% =====================================================================
%  D1 · PROJECT PROPOSAL — Big Data Analytics · Course Project
%  Project 7 · Customer value and baskets · Dataset: Online Retail II (UCI)
% =====================================================================
%
%  INSTRUCCIONES GENERALES PARA LOS AGENTES
%  - Fuente de verdad: ../project_guide.pdf (Secciones 3, 4, 6 y 7).
%  - El documento cubre EXACTAMENTE los 6 puntos del checklist D1 de la guía;
%    cada sección de abajo corresponde a uno de ellos.
%  - Idioma del texto final: inglés. Estas instrucciones están en español.
%  - Límite: máximo 2 páginas en total (portada incluida). Compilar con
%    `latexmk -pdf main.tex` y comprobar el número de páginas.
%  - Rúbrica D1 (5 puntos): "Clear questions and KPIs; suitable data (size,
%    license); realistic plan".
%  - No inventar datos personales (nombres, IDs, reparto de tareas): pedirlos
%    al usuario. Verificar cualquier cifra del dataset antes de escribirla.
%  - Sustituir cada \placeholder{...} por el contenido final.
% =====================================================================

\documentclass[11pt,a4paper]{article}

\usepackage[utf8]{inputenc}
\usepackage[T1]{fontenc}
\usepackage{lmodern}
\usepackage[english]{babel}
\usepackage[a4paper,margin=2cm]{geometry}
\usepackage{booktabs}
\usepackage{tabularx}
\usepackage{enumitem}
\usepackage{xcolor}
\usepackage{titlesec}
\usepackage[hidelinks]{hyperref}

\setlist{nosep, leftmargin=*}
\setlength{\parindent}{0pt}
\setlength{\parskip}{4pt}
\titlespacing*{\section}{0pt}{8pt}{3pt}
\titleformat{\section}{\large\bfseries}{\thesection.}{0.5em}{}

% Contenido pendiente: se imprime en rojo y debe desaparecer en la versión final.
\newcommand{\placeholder}[1]{\textcolor{red}{\textit{[TODO: #1]}}}

\begin{document}

% =====================================================================
%  1. PORTADA — Pair names and IDs; project number or own dataset
% =====================================================================
%  Qué debe ir aquí:
%  - Bloque compacto al inicio de la página 1 (no una página completa, por
%    el límite de 2 páginas).
%  - Título del proyecto, nombres completos e IDs de los dos estudiantes,
%    "Project 7 · Customer value and baskets", dataset Online Retail II (UCI)
%    y fecha de entrega (04-11-2026).
% =====================================================================
%  Estado: COMPLETADA, salvo \placeholder del segundo estudiante y de la
%  universidad. No cambiar el título sin consultar al usuario (eligió un
%  título académico, sin juegos de palabras).
\begin{center}
  {\small Universitat Rovira i Virgili \,·\, Big Data Analytics \,·\, Course Project}\\[4pt]
  {\Large\bfseries D1 · Project Proposal}\\[6pt]
  {\large\bfseries Identifying High-Value Customers, Product Associations and\\
  Future Spending in Online Retail Using Apache Spark}\\[8pt]
  \begin{tabular}{@{}ll@{\hspace{2.5em}}ll@{}}
    \toprule
    \textbf{Student} & \textbf{ID} & \textbf{Project} & 7 · Customer value and baskets \\
    \midrule
    Omar Morales Ortega & Y7350495P & \textbf{Dataset} & Online Retail II (UCI) \\
    \placeholder{Partner full name} & \placeholder{ID} \\
    \bottomrule
  \end{tabular}
\end{center}

% =====================================================================
\section{Business problem}
% =====================================================================
%  Qué debe ir aquí (checklist punto 2):
%  - Quién usará los resultados (un rol concreto dentro de la empresa).
%  - Para qué decisión o decisiones los usará.
%  - Breve contexto del negocio (tienda online de regalos) y por qué el
%    problema es relevante.
% =====================================================================
\placeholder{Who will use the results and for what decision.}

% =====================================================================
\section{Analysis, prediction and KPIs}
% =====================================================================
%  Qué debe ir aquí (checklist punto 3):
%  - Qué se va a analizar: las preguntas de análisis, alineadas con el
%    proyecto 7 (mejores clientes por recencia/frecuencia/gasto y productos
%    que se compran juntos). Incluir la pregunta "why?" que exige la guía.
%  - Qué se va a predecir: el gasto de cada cliente en el próximo trimestre
%    (lo que pide el proyecto 7), con su baseline.
%  - Los KPIs (mínimo 3 según la guía), cada uno con una definición clara.
% =====================================================================
\placeholder{What you will analyze, what you will predict, and the KPIs.}

% =====================================================================
\section{Data and platform}
% =====================================================================
%  Qué debe ir aquí (checklist punto 4):
%  - Fuente, link, licencia y tamaño del dataset Online Retail II (UCI).
%  - Confirmar que cumple el requisito de la guía (≥ 1 millón de filas de
%    datos reales, uso legal).
%  - Dónde se ejecutará: Google Colab o Databricks Free Edition (con PySpark).
% =====================================================================
\placeholder{Source, link, license, size; where you will run it (Colab or Databricks).}

% =====================================================================
\section{Advanced option and main risks}
% =====================================================================
%  Qué debe ir aquí (checklist punto 5):
%  - Opción avanzada elegida (A · Machine learning, la del proyecto 7) y el
%    plan para llevarla a cabo, incluyendo cómo se comparará con un baseline.
%  - Los riesgos principales del proyecto y cómo se mitigarán.
% =====================================================================
\placeholder{Advanced option (A--E) and your plan; main risks.}

% =====================================================================
\section{Team roles}
% =====================================================================
%  Qué debe ir aquí (checklist punto 6):
%  - Quién lidera cada parte del proyecto.
%  - Dejar claro que ambos miembros programan.
% =====================================================================
\placeholder{Who leads which part (both partners code).}

\end{document}
